\section{Inline Visual Route: Domain and Systems Readings}

These figures distribute domain examples across the manuscript instead of isolating them in an atlas.


\subsection{Physics replay lane}

This visual anchor names the objects, review direction, and formula or numeric boundary that keep the figure tied to the argument.


\begin{figure}[htbp]

\centering

\begin{tikzpicture}[>=Latex,node distance=2.4cm]

\node[draw,rounded corners,fill=blue!7,text width=0.28\textwidth,align=center,minimum height=1.0cm] (a) {Constant route};

\node[draw,rounded corners,fill=green!7,text width=0.28\textwidth,align=center,minimum height=1.0cm,right=of a] (b) {Comparator};

\draw[->,line width=0.8pt] (a) -- node[above,align=center] {residual} (b);

\node[draw,rounded corners,fill=orange!8,text width=0.68\textwidth,align=center,below=0.9cm of $(a)!0.5!(b)$] (c) {review question: support class, boundary, falsifier};

\draw[->,dashed] (c.north) -- ($(a)!0.5!(b)$);

\end{tikzpicture}

\caption{Physics replay lane. Shows bounded physical reconstruction as a replay lane, not total physics. The figure is placed inline in the manuscript; complete machine evidence remains in the evidence package.}

\label{fig:r005-inline-28d_oc133_inline_figures_domain_route-1}

\end{figure}

\subsection{Chemistry replay lane}

This visual anchor names the objects, review direction, and formula or numeric boundary that keep the figure tied to the argument.


\begin{figure}[htbp]

\centering

\begin{tikzpicture}[>=Latex,node distance=2.4cm]

\node[draw,rounded corners,fill=blue!7,text width=0.28\textwidth,align=center,minimum height=1.0cm] (a) {Formula mass};

\node[draw,rounded corners,fill=green!7,text width=0.28\textwidth,align=center,minimum height=1.0cm,right=of a] (b) {Public snapshot};

\draw[->,line width=0.8pt] (a) -- node[above,align=center] {residual} (b);

\node[draw,rounded corners,fill=orange!8,text width=0.68\textwidth,align=center,below=0.9cm of $(a)!0.5!(b)$] (c) {review question: support class, boundary, falsifier};

\draw[->,dashed] (c.north) -- ($(a)!0.5!(b)$);

\end{tikzpicture}

\caption{Chemistry replay lane. Shows molecular reconstruction with source and comparator. The figure is placed inline in the manuscript; complete machine evidence remains in the evidence package.}

\label{fig:r005-inline-28d_oc133_inline_figures_domain_route-2}

\end{figure}

\subsection{Biology replay lane}

This visual anchor names the objects, review direction, and formula or numeric boundary that keep the figure tied to the argument.


\begin{figure}[htbp]

\centering

\begin{tikzpicture}[>=Latex,node distance=2.4cm]

\node[draw,rounded corners,fill=blue!7,text width=0.28\textwidth,align=center,minimum height=1.0cm] (a) {Observed count};

\node[draw,rounded corners,fill=green!7,text width=0.28\textwidth,align=center,minimum height=1.0cm,right=of a] (b) {Protocol};

\draw[->,line width=0.8pt] (a) -- node[above,align=center] {negative control} (b);

\node[draw,rounded corners,fill=orange!8,text width=0.68\textwidth,align=center,below=0.9cm of $(a)!0.5!(b)$] (c) {review question: support class, boundary, falsifier};

\draw[->,dashed] (c.north) -- ($(a)!0.5!(b)$);

\end{tikzpicture}

\caption{Biology replay lane. Shows bio/geodata evidence as bounded QA. The figure is placed inline in the manuscript; complete machine evidence remains in the evidence package.}

\label{fig:r005-inline-28d_oc133_inline_figures_domain_route-3}

\end{figure}

\subsection{Systems replay lane}

This visual anchor names the objects, review direction, and formula or numeric boundary that keep the figure tied to the argument.


\begin{figure}[htbp]

\centering

\begin{tikzpicture}[>=Latex,node distance=2.4cm]

\node[draw,rounded corners,fill=blue!7,text width=0.28\textwidth,align=center,minimum height=1.0cm] (a) {Macro series};

\node[draw,rounded corners,fill=green!7,text width=0.28\textwidth,align=center,minimum height=1.0cm,right=of a] (b) {Baseline};

\draw[->,line width=0.8pt] (a) -- node[above,align=center] {forecast boundary} (b);

\node[draw,rounded corners,fill=orange!8,text width=0.68\textwidth,align=center,below=0.9cm of $(a)!0.5!(b)$] (c) {review question: support class, boundary, falsifier};

\draw[->,dashed] (c.north) -- ($(a)!0.5!(b)$);

\end{tikzpicture}

\caption{Systems replay lane. Shows systems evidence as protocol-bound series reading. The figure is placed inline in the manuscript; complete machine evidence remains in the evidence package.}

\label{fig:r005-inline-28d_oc133_inline_figures_domain_route-4}

\end{figure}

\subsection{AI builder route}

This visual anchor names the objects, review direction, and formula or numeric boundary that keep the figure tied to the argument.


\begin{figure}[htbp]

\centering

\begin{tikzpicture}[>=Latex,node distance=2.4cm]

\node[draw,rounded corners,fill=blue!7,text width=0.28\textwidth,align=center,minimum height=1.0cm] (a) {Model grammar};

\node[draw,rounded corners,fill=green!7,text width=0.28\textwidth,align=center,minimum height=1.0cm,right=of a] (b) {Evaluation pipeline};

\draw[->,line width=0.8pt] (a) -- node[above,align=center] {claim boundary} (b);

\node[draw,rounded corners,fill=orange!8,text width=0.68\textwidth,align=center,below=0.9cm of $(a)!0.5!(b)$] (c) {review question: support class, boundary, falsifier};

\draw[->,dashed] (c.north) -- ($(a)!0.5!(b)$);

\end{tikzpicture}

\caption{AI builder route. Shows how engineers use OC as claim-to-evidence discipline. The figure is placed inline in the manuscript; complete machine evidence remains in the evidence package.}

\label{fig:r005-inline-28d_oc133_inline_figures_domain_route-5}

\end{figure}

\subsection{Enterprise route}

This visual anchor names the objects, review direction, and formula or numeric boundary that keep the figure tied to the argument.


\begin{figure}[htbp]

\centering

\begin{tikzpicture}[>=Latex,node distance=2.4cm]

\node[draw,rounded corners,fill=blue!7,text width=0.28\textwidth,align=center,minimum height=1.0cm] (a) {Decision question};

\node[draw,rounded corners,fill=green!7,text width=0.28\textwidth,align=center,minimum height=1.0cm,right=of a] (b) {Evidence class};

\draw[->,line width=0.8pt] (a) -- node[above,align=center] {commitment boundary} (b);

\node[draw,rounded corners,fill=orange!8,text width=0.68\textwidth,align=center,below=0.9cm of $(a)!0.5!(b)$] (c) {review question: support class, boundary, falsifier};

\draw[->,dashed] (c.north) -- ($(a)!0.5!(b)$);

\end{tikzpicture}

\caption{Enterprise route. Shows how CIO/CEO/EA readers avoid overcommitment. The figure is placed inline in the manuscript; complete machine evidence remains in the evidence package.}

\label{fig:r005-inline-28d_oc133_inline_figures_domain_route-6}

\end{figure}
